\documentclass[a4paper,11pt]{article}
\usepackage{exam-style-842-unit}
\examinehead{842 信息化概论 · 第 2 单元练习册}{练习二\quad 信息系统与数据资源管理}
\begin{document}

\unitcover{842 信息化概论 · 单元练习册}{第 2 单元\quad 数据管理、数据库与信息系统}{练习二\quad 信息系统与数据资源管理（共 100 分）}{约 70 分钟}{信息系统的概念与组成要素；信息系统的五个基本功能；信息系统的类型（TPS、MIS、DSS）；信息资源与信息资源管理的含义、目标与内容；信息资源的获取、加工、存储、检索、传输与共享；数据管理与数据库管理；信息安全与数据保护；数据备份与信息孤岛。}

\begin{knowmap}
\textbf{信息系统：}由人、硬件、软件、数据和处理规程五个要素组成，对信息进行采集、加工、存储、传输和利用的人机系统\\[2pt]
\textbf{五个基本功能：}输入 $\rightarrow$ 处理 $\rightarrow$ 输出，并由存储与控制贯穿其中\\[2pt]
\textbf{按功能与层次分：}事务处理系统（TPS）$\rightarrow$ 管理信息系统（MIS）$\rightarrow$ 决策支持系统（DSS）$\rightarrow$ 专家系统（ES）\\[2pt]
\textbf{信息资源管理：}把信息视为与物质、能源同等重要的资源，对信息、信息人员、信息技术与设施等要素进行计划、组织、协调和控制\\[2pt]
\textbf{信息资源管理环节：}采集 $\rightarrow$ 加工 $\rightarrow$ 存储 $\rightarrow$ 检索 $\rightarrow$ 传输 $\rightarrow$ 共享利用，全程贯穿安全保护\\[2pt]
\textbf{数据保护：}机密性 · 完整性 · 可用性；手段包括加密、认证、访问控制、备份、审计与管理制度
\end{knowmap}

\begin{keyterms}
\textbf{信息系统}：由人、硬件、软件、数据和处理规程组成，对信息进行采集、加工、存储、传输和利用的人机系统。\\
\textbf{信息资源}：人类社会信息活动中积累起来的信息本身以及信息生产者、信息技术与信息设施等信息活动要素的集合。\\
\textbf{信息资源管理（IRM）}：对信息活动中的信息、人员、技术、设备和资金等要素进行计划、组织、协调与控制，以实现信息资源的合理开发与有效利用。\\
\textbf{信息采集}：通过传感、调查、统计、录入等手段获取原始信息的过程，要求全面、准确、及时。\\
\textbf{信息加工}：对采集到的原始信息进行分类、排序、计算、分析、综合等处理，使之成为对使用者有用的信息。\\
\textbf{信息存储}：把信息以一定格式记录在相应载体上长期保存，以便反复使用，要求安全、可靠、便于检索。\\
\textbf{信息检索}：从信息集合中查找所需信息的过程，有目录检索、关键词检索、全文检索等方式。\\
\textbf{信息共享}：信息可被多个用户同时使用而不损耗，共享范围越大，信息的价值越高。\\
\textbf{数据库管理员（DBA）}：负责数据库的建立、维护、监控、备份与安全管理，并协调数据库使用的人员。\\
\textbf{数据备份}：为防止数据丢失或损坏，把数据复制到其他存储介质上保存的保护措施。\\
\textbf{数据安全}：保护数据不被非法获取、篡改、泄露或破坏，包括机密性、完整性和可用性等基本属性。\\
\textbf{事务处理系统（TPS）}：面向组织日常业务、处理大量结构化事务数据的信息系统，如订单、账务、库存系统。\\
\textbf{管理信息系统（MIS）}：在事务处理系统基础上对管理信息进行加工，为各级管理者提供报表、查询和管理支持的系统。\\
\textbf{决策支持系统（DSS）}：以数据和模型为手段，辅助管理者进行半结构化决策的人机交互系统。
\end{keyterms}

\begin{contrast}{事务处理系统 / 管理信息系统 / 决策支持系统（每行依次为：处理对象·主要功能·服务对象）}
\citem{事务处理系统（TPS）}{日常业务中的大量结构化事务数据}{数据的录入、更新、查询与统计，保证业务实时处理}{基层业务人员}
\citem{管理信息系统（MIS）}{由事务数据汇总而成的管理信息}{生成报表、进行查询与例行管理控制}{中层管理者}
\citem{决策支持系统（DSS）}{数据与模型、方法库}{以人机交互方式支持半结构化、非结构化决策}{高层决策者}
\end{contrast}

\begin{contrast}{数据管理 / 信息资源管理（每行依次为：管理对象·主要目标·主要手段）}
\citem{数据管理}{数据库或数据文件中的数据}{保证数据可共享、冗余度小、独立性强、安全可靠}{数据库管理系统、数据库设计与 SQL、备份与恢复}
\citem{信息资源管理}{信息内容、信息人员、信息技术与信息设施}{把信息作为战略资源进行合理配置和有效利用，提高组织效益}{管理体制、政策法规、标准规范与信息技术相结合}
\end{contrast}

\begin{contrast}{数据完整性 / 数据安全（每行依次为：关注重点·主要威胁·典型措施）}
\citem{数据完整性}{数据正确、一致、有效}{错误数据、非法数据进入数据库}{实体完整性、参照完整性、用户定义完整性约束}
\citem{数据安全}{数据不被非法获取、篡改、泄露和破坏}{非授权访问、篡改、泄露、病毒与灾害}{加密、认证、访问控制、备份、审计与安全管理制度}
\end{contrast}

\begin{contrast}{信息采集 / 信息加工 / 信息存储（信息资源管理的主要环节；每行依次为：环节含义·主要方法·基本要求）}
\citem{信息采集}{获取原始信息}{传感采集、调查统计、数据录入}{全面、准确、及时}
\citem{信息加工}{把原始信息变成有用信息}{分类、排序、计算、分析、综合}{真实、适用、及时}
\citem{信息存储}{把信息长期保存以备再用}{文件、数据库、磁盘与光盘、云存储}{安全、可靠、便于检索}
\end{contrast}

\qtypename{一、单项选择题}{10 题，每题 3 分，共 30 分}
\q 下列不属于信息系统五个基本功能的是（\quad）
\choicesline{A. 信息的输入}{B. 信息的处理}{C. 信息的存储}{D. 资金的筹措}

\q 信息系统由人、硬件、软件、数据和（\quad）五个要素组成。
\choicesline{A. 网络}{B. 处理规程}{C. 机房}{D. 报表}

\q 信息资源管理的核心是（\quad）
\choices{A. 把所有纸质文件数字化}{B. 把信息作为资源进行合理开发与有效利用}
{C. 只购买先进的计算机设备}{D. 只建设数据库}

\q 信息采集环节的基本要求是（\quad）
\choicesline{A. 全面、准确、及时}{B. 大量、快速、廉价}{C. 隐蔽、安全、保密}{D. 简单、随意、方便}

\q 面向组织日常业务、处理大量结构化事务数据的信息系统是（\quad）
\choicesline{A. 事务处理系统}{B. 管理信息系统}{C. 决策支持系统}{D. 专家系统}

\q 在事务处理系统基础上对管理信息进行加工，为各级管理者提供报表和管理支持的系统是（\quad）
\choicesline{A. 事务处理系统}{B. 管理信息系统}{C. 决策支持系统}{D. 办公自动化系统}

\q 以数据和模型为手段、辅助管理者进行半结构化决策的系统是（\quad）
\choicesline{A. 事务处理系统}{B. 管理信息系统}{C. 决策支持系统}{D. 数据库系统}

\q 信息安全的基本属性是（\quad）
\choices{A. 机密性、完整性、可用性}{B. 开放性、共享性、透明性}
{C. 全面性、准确性、及时性}{D. 经济性、方便性、灵活性}

\q 为防止数据丢失或损坏，把数据复制到其他存储介质上保存的措施称为（\quad）
\choicesline{A. 数据加密}{B. 数据备份}{C. 数据压缩}{D. 数据索引}

\q 同一信息可以被多个用户同时使用而不损耗，这体现了信息的（\quad）
\choicesline{A. 可存储性}{B. 可共享性}{C. 可压缩性}{D. 可度量性}

\qtypename{二、多项选择题}{3 题，每题 4 分，共 12 分；每题有两个或两个以上正确项}
\q 信息系统的组成要素包括（\quad）
\choicesline{A. 人}{B. 硬件与软件}{C. 数据}{D. 资金}

\q 信息资源管理的主要内容包括（\quad）
\choices{A. 信息资源的开发与利用}{B. 信息基础设施与信息系统的建设}
{C. 信息安全与保密管理}{D. 只管理纸质档案}

\q 信息安全的基本属性包括（\quad）
\choicesline{A. 机密性}{B. 完整性}{C. 可用性}{D. 随意性}

\qtypename{三、判断题}{6 题，每题 2 分，共 12 分；正确的画“√”，错误的画“×”并写出改正后的表述}
\q 信息系统就是计算机硬件系统。\quad（\quad）\quad 改正：\fixblank

\q 信息资源共享会因多人使用而使信息总量减少。\quad（\quad）\quad 改正：\fixblank

\q 数据备份是防止数据丢失和损坏的有效措施。\quad（\quad）\quad 改正：\fixblank

\q 信息系统的输出只能是打印出来的纸质报表。\quad（\quad）\quad 改正：\fixblank

\q 信息的采集应做到全面、准确、及时。\quad（\quad）\quad 改正：\fixblank

\q 决策支持系统以数据和模型为手段，辅助管理者进行半结构化决策。\quad（\quad）\quad 改正：\fixblank

\qtypename{四、填空题}{6 空，每空 1 分，共 6 分}
\q 信息系统的五个基本功能是输入、处理、输出、存储和\fillblanks[2.2cm]。

\q 信息系统由人、硬件、软件、数据和\fillblanks[2.4cm]五个要素组成。

\q 信息资源的开发利用包括信息的采集、加工、存储、检索、传输和\fillblanks[1.6cm]。

\q 信息安全的基本属性是机密性、完整性和\fillblanks[1.8cm]。

\q 面向组织日常业务、主要处理结构化事务数据的信息系统称为\fillblanks[2.6cm]。

\q 把数据复制到其他存储介质上保存，以防止数据丢失的保护措施称为\fillblanks[2.2cm]。

\newpage
\qtypename{五、简答题}{4 题，每题 6 分，共 24 分}
\q 什么是信息系统？简述信息系统的组成要素与五个基本功能。\anslines{2.4cm}

\q 什么是信息资源管理？简述信息资源管理的主要目标与内容。\anslines{2.4cm}

\q 简述信息资源管理的六个主要环节及其基本要求。\anslines{2.4cm}

\q 信息安全的基本属性有哪些？常用的数据保护措施有哪些？\anslines{2.4cm}

\newpage
\qtypename{六、综合题（论述）}{1 题，共 16 分}
\q 某县农业农村局建设“农业信息资源平台”，把土壤、气象、市场行情、种植户等数据集中存储，向乡镇农技站、合作社、龙头企业开放共享，同时制定数据安全管理制度。请结合这一案例回答：
（1）说明该平台作为信息系统，其组成要素和所实现的基本功能；
（2）说明信息资源管理的目标，以及该平台应包含哪些主要环节；
（3）分析在开放共享的同时如何保障数据安全；
（4）指出该平台建设中可能出现的信息孤岛问题，并提出解决对策。\anslines{5cm}

\end{document}
